\documentclass[11pt,a4paper]{article}
\usepackage[margin=25mm]{geometry}
\usepackage[T1]{fontenc}
\usepackage{lmodern,microtype,amsmath,amssymb,amsthm,booktabs,array,longtable}
\usepackage[dvipsnames]{xcolor}
\usepackage[colorlinks=true,linkcolor=MidnightBlue,citecolor=MidnightBlue,urlcolor=MidnightBlue]{hyperref}
\usepackage{enumitem}
\setlist{itemsep=2pt,topsep=5pt}
\newtheorem{theorem}{Theorem}
\newtheorem{lemma}{Lemma}
\newtheorem{remark}{Remark}
\newcommand{\F}{\mathbb F_2}
\newcommand{\wt}{\operatorname{wt}}
\newcommand{\rank}{\operatorname{rank}}
\newcommand{\Tr}{\operatorname{Tr}}
\newcommand{\rad}{\operatorname{rad}}
\newcommand{\Span}{\operatorname{span}_{\F}}
\newcommand{\CA}{\mathcal Q_A}
\newcommand{\CB}{\mathcal Q_B}
\newcommand{\CC}{\mathcal Q_C}
\newcommand{\pauli}[1]{\texttt{#1}}
\title{Three locally inequivalent pure $[[18,3,6]]$ stabilizer codes\\[3pt]\large Explicit gluing, signed Clifford transformations, and exact certificates}
\author{Brandon Li}
\date{}
\begin{document}
\maketitle
\vspace{-1.5em}
\begin{abstract}
We give three explicit pure binary $[[18,3,6]]$ stabilizer codes and prove that they are pairwise inequivalent under products of single-qubit unitaries and qubit permutations. The first construction glues eight- and ten-qubit binary symplectic spaces through ten-dimensional quotients whose dual alternating forms have rank eight. Two weight-four symplectic transvections produce the other codes; explicit Pauli corrections fix every generator sign in the corresponding Clifford conjugations. The stabilizer minimum weights are $6,7,7$, and the distinguishing stabilizer coefficients are $a_6=4$, $a_7=16$, and $a_7=12$, respectively. We provide all generators, complete stabilizer and normalizer weight enumerators, weight-six logical witnesses, and a portable exact-arithmetic certificate. The distance proofs combine a transparent local coset bound with independent exhaustive checks in the physical Pauli space. The contribution is the explicit, mutually locally inequivalent constructions and their reproducible certificates.
\end{abstract}

\section{Introduction}
Small stabilizer codes are useful test cases for constructive methods, equivalence invariants, and exact verification. At short lengths it is feasible to expose every generator and check the full normalizer rather than infer a distance from a heuristic search. This makes it possible to separate three questions: whether a particular code has the claimed parameters, how it is constructed, and whether two presentations describe locally equivalent code spaces.

This paper answers those questions for three binary codes of length $18$, dimension $2^3$, and distance $6$. Code $\CA$ arises from a local-to-global gluing construction. Codes $\CB$ and $\CC$ are its images under explicitly specified, nonlocal Clifford operations. Despite this global Clifford relationship, the three codes are inequivalent under arbitrary products of single-qubit unitaries, even when permutations of the qubits are allowed. Complete weight enumerators give short witnesses to this inequivalence.

The underlying stabilizer formalism and its description by binary symplectic or additive $\mathrm{GF}(4)$ codes are standard \cite{gottesman,calderbank}. Quantum weight enumerators and their MacWilliams relations provide the consistency checks and equivalence invariant used here \cite{shor-laflamme,rains}. The gluing presentation below is stated directly in binary coordinates, including the distinction between quotient spaces and their dual alternating forms.

\paragraph{Prior bounds and prior claims.}
The standard additive-code bounds of Calderbank, Rains, Shor, and Sloane list $5\le d\le6$ at $(n,k)=(18,3)$; in the inspected preprint version, Table III and its notes trace the unmarked upper bound to the linear-programming constraints of Theorem 21 \cite{calderbank}. The three codes here attain that established upper bound for binary stabilizer codes. Grassl's online entry also displayed lower bound $5$ and upper bound $6$ when accessed on 11 November 2026 \cite{GrasslTable183}. Thus the explicit certificates improve the lower bound displayed in that entry.

In particular, Cohen, Encheva, and Litsyn already explicitly claimed $[[18,3,6]]$ in Example 1 of their 1998 preprint \cite{CohenEnchevaLitsyn1998preprint}. The same claim and displayed matrix appear on page 2496 of the final 1999 article \cite{CohenEnchevaLitsyn1999}. Direct inspection and exact computation show that the printed $12$-dimensional extension has Euclidean Gram rank $10$ and cannot contain the required $9$-dimensional dual-containing subcode. Thus those displayed data do not satisfy the invoked theorem's hypothesis. Appendix~\ref{app:prior} gives the precise obstruction. This paper gives three explicit constructions, their local inequivalence, and their complete certificates.

The manuscript contains three specific constructions, not a classification of length-$18$ codes. In particular, it neither constructs a $[[18,4,6]]$ code nor rules one out. A global Clifford map preserves commutation and dimension but need not preserve physical Pauli weight, so the distance of each derivative is verified separately.

\paragraph{Organization.}
Section~\ref{sec:notation} fixes conventions. Section~\ref{sec:gluing} gives the gluing lemma and the concrete construction of $\CA$. Section~\ref{sec:derivatives} defines the other two codes with exact phases. Section~\ref{sec:certificates} proves the parameters and records all enumerators. Section~\ref{sec:inequivalence} proves local-unitary inequivalence. Section~\ref{sec:reproducibility} describes the finite certificates and their reproduction. The appendices display every stabilizer generator and the full local gluing data.

\section{Binary symplectic and Pauli conventions}\label{sec:notation}
Let $V_n=\F^n\oplus\F^n$. Write $r=(x\mid z)$ and use
\begin{equation}\label{eq:pair}
 [r,r']=x\cdot z'+z\cdot x',\qquad
 \wt(r)=\#\{i:(x_i,z_i)\ne(0,0)\}.
\end{equation}
Coordinates in the paper are numbered $1,\ldots,n$. The corresponding integer encoding is
\begin{equation}\label{eq:int}
 \operatorname{enc}_n(r)=\sum_{i=1}^n x_i2^{i-1}+\sum_{i=1}^n z_i2^{n+i-1}.
\end{equation}
Thus the leftmost character of a displayed Pauli word is the least significant coordinate of each binary half. The pairs $(0,0),(1,0),(0,1),(1,1)$ mean $I,X,Z,Y$. All additions of binary vectors, including encoded integers, are XOR operations.

We associate to $r$ the positive Hermitian tensor product
\begin{equation}\label{eq:hermitian}
 P_r=i^{q(r)}X^xZ^z,\qquad q(r)=\sum_{i=1}^n x_i z_i\in\mathbb Z.
\end{equation}
Here $q(r)$ is evaluated as an ordinary integer when used in a phase. Then $P_rP_{r'}=(-1)^{[r,r']}P_{r'}P_r$.

For an isotropic binary subspace $S\subseteq V_n$ of rank $n-k$, select a binary basis $s_1,\ldots,s_{n-k}$ and impose the $+1$ eigenvalue of each $P_{s_j}$. These independent commuting Hermitian operators generate an abelian Pauli group $\mathcal S$ not containing $-I$. Indeed, a nonempty product with scalar binary label would give a nontrivial dependence among the $s_j$. The joint eigenspace has dimension $K=2^k$. The binary normalizer is $N=S^\perp$, and
\begin{equation}\label{eq:distance}
 d=\min\{\wt(r):r\in N\setminus S\}.
\end{equation}
We call the code pure to distance $d$ if no nonzero stabilizer label has weight less than $d$. In particular, $\min_{0\ne r\in N}\wt(r)=d$ together with a weight-$d$ member of $N\setminus S$ proves purity and exact distance.

For a code $\mathcal Q$ define the integer counts
\begin{equation}\label{eq:enumerators}
 a_j(\mathcal Q)=\#\{s\in S:\wt(s)=j\},\quad
 b_j(\mathcal Q)=\#\{r\in N:\wt(r)=j\},\quad
 \ell_j(\mathcal Q)=b_j(\mathcal Q)-a_j(\mathcal Q).
\end{equation}
The last quantity counts binary Pauli labels in $N\setminus S$, not logical equivalence classes modulo $S$. The associated homogeneous weight enumerators are $A(u,v)=\sum_j a_ju^{n-j}v^j$ and $B(u,v)=\sum_j b_ju^{n-j}v^j$.

\section{An eight-plus-ten gluing construction}\label{sec:gluing}
\subsection{Quotients and their dual alternating forms}
Take subspaces $C\subseteq V_{n_L}$ and $E\subseteq V_{n_R}$, and marked spaces
\[
 A\subseteq C\cap C^\perp,\qquad B\subseteq E\cap E^\perp.
\]
The letters $A,B$ in this section denote marked binary subspaces, not code labels or enumerator polynomials. Suppose
\begin{equation}\label{eq:quotients}
 W_L=A^\perp/C,\qquad W_R=B^\perp/E
\end{equation}
have a common dimension $m$. Choose linear lift maps $q_L:\F^m\to A^\perp$ and $q_R:\F^m\to B^\perp$ whose cosets give quotient bases.

The ambient pairing need not descend to the primal quotients in~\eqref{eq:quotients}. Instead, the natural dual identifications are
\begin{equation}\label{eq:dualquotients}
 W_L^*\cong C^\perp/A,\qquad W_R^*\cong E^\perp/B.
\end{equation}
For example, $d\in C^\perp$ defines the functional $t\mapsto[d,q_L(t)]$, and its kernel as a map from $C^\perp$ to $W_L^*$ is exactly $A$. Since $A$ is in the radical of the restricted form on $C^\perp$, this dual quotient does carry an alternating form. Let $J_L,J_R$ be its matrices in the coordinate dual bases. Equivalently, choose $d_{L,i}\in C^\perp$ with $[d_{L,i},q_L(e_j)]=\delta_{ij}$, and set $(J_L)_{ij}=[d_{L,i},d_{L,j}]$; define $J_R$ similarly. These are the quotient Poisson matrices. Their definition is independent of the chosen representatives modulo $A$ and $B$.

\begin{lemma}[Coisotropic gluing]\label{lem:gluing}
Let $F:\F^m\to\F^m$ be invertible and satisfy
\begin{equation}\label{eq:poisson}
 FJ_LF^{\mathsf T}=J_R.
\end{equation}
In $V_{n_L}\oplus V_{n_R}$ define
\begin{equation}\label{eq:globalN}
 N=(C\oplus E)+\{q_L(t)\oplus q_R(Ft):t\in\F^m\}.
\end{equation}
Then $\dim N=\dim C+\dim E+m$, and $S=N^\perp$ is isotropic. Hence $S\subseteq N=S^\perp$ defines a stabilizer code.

For $t\in\F^m$ put
\[
 c_L(t)=\min_{c\in C}\wt(q_L(t)+c),\qquad
 c_R(t)=\min_{e\in E}\wt(q_R(t)+e).
\]
If every nonzero member of $C$ and $E$ has weight at least $d_0$, and
\begin{equation}\label{eq:cost}
 c_L(t)+c_R(Ft)\ge d_0\quad(t\ne0),
\end{equation}
then every nonzero member of $N$ has weight at least $d_0$.
\end{lemma}
\begin{proof}
The $m$ paired quotient lifts are independent modulo $C\oplus E$, proving the dimension formula. A vector in $S$ has components in $C^\perp$ and $E^\perp$. Let their coordinate functionals on the two quotients be $\alpha$ and $\beta$. Orthogonality to every graph lift gives
\[
 \alpha^{\mathsf T}t+\beta^{\mathsf T}Ft=0\quad\text{for all }t,
 \qquad\text{so}\qquad \alpha=F^{\mathsf T}\beta.
\]
For two such vectors, the ambient symplectic pairing equals
\[
 \alpha^{\mathsf T}J_L\alpha'+\beta^{\mathsf T}J_R\beta'
 =\beta^{\mathsf T}(FJ_LF^{\mathsf T}+J_R)\beta'=0.
\]
The terms from $A$ and $B$ vanish by their radical property. Thus $S$ is isotropic. Nondegeneracy of the ambient space gives $S^\perp=N$.

Every vector of $N$ has a unique quotient label $t$. For $t\ne0$, the disjoint coordinate blocks make its weight at least $c_L(t)+c_R(Ft)$. If $t=0$, a nonzero vector is a member of $C\oplus E$ with at least one nonzero component, and the stated local minimum bound applies.
\end{proof}

\subsection{Concrete local data}
We take $(n_L,n_R)=(8,10)$ and use the local bases in Appendix~\ref{app:local}. The four displayed generators of $C$ are independent and pairwise orthogonal. The first two span $A$. The seven displayed generators of $E$ are independent and have Gram rank $2$; the first three span $B\subseteq\rad(E)$. Consequently,
\[
 \dim(A^\perp/C)=16-2-4=10,
 \qquad \dim(B^\perp/E)=20-3-7=10.
\]
The dimensions of the dual radicals are both $2$: on the left they are those of $C/A$, and on the right those of $(E\cap E^\perp)/B$, since $\dim(E\cap E^\perp)=7-2=5$. The forms on the ten-dimensional dual quotient spaces therefore have rank $8$.

For the stated quotient lifts, the exact dual matrices coincide:
\begin{equation}\label{eq:J}
 J_L=J_R=J=\begin{pmatrix}0&1\\1&0\end{pmatrix}^{\oplus4}\oplus 0_2.
\end{equation}
Dual representatives certifying this assertion are included in Appendix~\ref{app:local}. The quotient map $F$ is specified by its ten integer-encoded columns:
\begin{equation}\label{eq:Fcols}
 (Fe_1,\ldots,Fe_{10})=(240,31,248,100,144,202,242,97,389,678).
\end{equation}
The columns have rank $10$ and satisfy $FJF^{\mathsf T}=J$ over $\F$. Their encodings use ten coordinate bits, with $e_1$ the low bit.

The local nonzero minimum weights are both $6$. In univariate notation their exact weight enumerators are
\begin{align}
 W_C(y)&=1+12y^6+3y^8,\\
 W_E(y)&=1+38y^6+16y^7+47y^8+16y^9+10y^{10}.
\end{align}
Every cost in~\eqref{eq:cost} is computed by minimizing over $16$ or $128$ local words. The complete distribution of the $1023$ nonzero quotient costs is
\begin{center}
\begin{tabular}{@{}lrrrrr@{}}\toprule
$c_L(t)+c_R(Ft)$&6&7&8&9&10\\
Number of nonzero $t$&261&333&337&73&19\\\bottomrule
\end{tabular}
\end{center}
Thus every nonzero vector of the resulting $N$ has weight at least $6$.

\begin{theorem}\label{thm:A}
The gluing data in Appendix~\ref{app:local} and~\eqref{eq:Fcols} define a pure binary $[[18,3,6]]$ stabilizer code $\CA$. A positive Hermitian generator basis is given in Appendix~\ref{app:generators}. Its nonzero stabilizer minimum weight is $6$, with $a_6(\CA)=4$, and $b_6(\CA)=588$.
\end{theorem}
\begin{proof}
Lemma~\ref{lem:gluing} gives $\dim N=4+7+10=21$ and an isotropic orthogonal complement $S=N^\perp$ of dimension $36-21=15$. Hence the joint eigenspace encodes $18-15=3$ qubits. The local cost bound proves that $N$ contains no nonzero vector of weight below $6$. Taking the symplectic orthogonal complement of the explicit $21$ spanning vectors in~\eqref{eq:globalN} gives precisely the $15$ independent rows displayed for $\CA$.

The word
\begin{equation}\label{eq:witnessA}
 P_{r_A}=\pauli{XIIIIXIIIXIIXIIIXY},\qquad
 \operatorname{enc}_{18}(r_A)=34359939617,
\end{equation}
has weight $6$, is orthogonal to every displayed generator, and is not in their binary span. These claims follow by binary elimination and are checked by the certificate in Section~\ref{sec:reproducibility}. Equation~\eqref{eq:distance} therefore gives distance exactly $6$ and purity. Enumeration of the $2^{15}$ stabilizer labels and $2^{21}$ normalizer labels yields the coefficients in Table~\ref{tab:enumerators}, including $a_6=4$ and $b_6=588$.
\end{proof}

\section{Two exact Clifford derivatives}\label{sec:derivatives}
For a binary Pauli label $v$, the symplectic transvection
\begin{equation}\label{eq:transvection}
 T_v(r)=r+[r,v]v
\end{equation}
is an involution preserving~\eqref{eq:pair}. Define $S_B=T_{v_B}(S_A)$ and $S_C=T_{v_C}(S_A)$ with
\begin{align}
 v_B&=4328808577,&P_{v_B}&=\pauli{YIIIIIIYIIIIIXYIII},\label{eq:vb}\\
 v_C&=8607273026,&P_{v_C}&=\pauli{IYIIIIYIIIIIXIIYII}.\label{eq:vc}
\end{align}
The nonidentity supports are $\{1,8,14,15\}$ and $\{2,7,13,16\}$, respectively. Both masks have weight four. Applying the transvections row by row to Appendix~\ref{app:generators} gives the displayed bases of $S_B$ and $S_C$, in the same order.

A binary symplectic map alone does not specify the signs of Hermitian generator images. For completeness, set
\begin{equation}\label{eq:unitary}
 U_v=\exp(i\pi P_v/4)=\frac{I+iP_v}{\sqrt2}.
\end{equation}
It implements $T_v$ on binary labels. When $[v,r]=0$, conjugation fixes $P_r$; otherwise it gives $iP_vP_r$, which is either $+P_{v+r}$ or $-P_{v+r}$.

Choose the correction labels
\begin{align}
 w_B&=17435,&P_{w_B}&=\pauli{XXIXXIIIIIXIIIXIII},\\
 w_C&=280636,&P_{w_C}&=\pauli{ZIXXXXIIIIIXIIXIII}.
\end{align}
For every row index $j=1,\ldots,15$, the exact, signed conjugations are
\begin{equation}\label{eq:exactclifford}
 (P_{w_D}U_{v_D})P_{s_{A,j}}(P_{w_D}U_{v_D})^\dagger
 =P_{s_{D,j}},\qquad D\in\{B,C\}.
\end{equation}
For $B$, bare $U_{v_B}$ introduces minus signs at rows $9,12$; for $C$, bare $U_{v_C}$ introduces them at rows $3,4,5$. Conjugation by the stated $P_{w_D}$ cancels exactly these signs. All row indices here are one-based.

The phase verification uses the integer-valued formula
\begin{equation}\label{eq:productphase}
 P_aP_b=i^{\phi(a,b)}P_{a+b},\quad
 \phi(a,b)=q(a)+q(b)-q(a+b)+2z_a\cdot x_b\pmod4.
\end{equation}
For an anticommuting row $s_{A,j}$ the bare sign is $(-1)^{(1+\phi(v_D,s_{A,j}))/2}$, where the even numerator is first reduced modulo four. The correction obeys
\[
 [w_D,s_{D,j}]=\frac{(1+\phi(v_D,s_{A,j}))\bmod4}{2}
\]
for these rows, and zero for the commuting rows. These are $30$ exact binary checks. The verifier additionally checks~\eqref{eq:productphase} against all $16$ one-qubit Pauli products as explicit $2\times2$ matrices.

\begin{theorem}\label{thm:BC}
The positive Hermitian generator lists in Appendix~\ref{app:generators} define pure binary $[[18,3,6]]$ codes $\CB$ and $\CC$. Their stabilizer minimum weights are both $7$, with
\[
 a_7(\CB)=16,\qquad a_7(\CC)=12,
 \qquad b_6(\CB)=b_6(\CC)=564.
\]
Equation~\eqref{eq:exactclifford} maps the code space of $\CA$ exactly onto the respective code spaces, including the chosen generator signs.
\end{theorem}
\begin{proof}
Symplectic transvections preserve rank and isotropy, so both stabilizers have rank $15$ and both codes encode three qubits. Equation~\eqref{eq:productphase} and the stated corrections prove~\eqref{eq:exactclifford}. Each derivative is then checked separately in physical coordinates: exhaustive enumeration of its $2^{21}$ normalizer labels gives zero at weights $1,\ldots,5$, while the explicit normalizer words outside the stabilizer are
\begin{align}
 \operatorname{enc}_{18}(r_B)&=34359943556,
 &P_{r_B}&=\pauli{IIXIIIIXXIIIIXIIXY},\label{eq:witnessB}\\
 \operatorname{enc}_{18}(r_C)&=34359939617,
 &P_{r_C}&=\pauli{XIIIIXIIIXIIXIIIXY}.\label{eq:witnessC}
\end{align}
Both have weight $6$. The complete enumerators in Table~\ref{tab:enumerators} show no nonzero stabilizer word below weight $7$ and give the stated coefficients. Thus the distance is exactly $6$ and both codes are pure. The independent low-weight error checks described below give a second proof of the absence of normalizer labels of weights at most five.
\end{proof}

\section{Complete physical weight enumerators}\label{sec:certificates}
Table~\ref{tab:enumerators} lists all stabilizer and normalizer coefficients, including every zero coefficient. The logical-label coefficients are exactly $\ell_j=b_j-a_j$, so the table also specifies the complete logical weight distributions without a further convention. The accompanying CSV additionally lists those differences explicitly.

\begin{table}[tbp]
\centering\small
\setlength{\tabcolsep}{6pt}
\begin{tabular}{@{}rrrrrrr@{}}\toprule
&\multicolumn{2}{c}{$\CA$}&\multicolumn{2}{c}{$\CB$}&\multicolumn{2}{c}{$\CC$}\\
\cmidrule(lr){2-3}\cmidrule(lr){4-5}\cmidrule(l){6-7}
$j$&$a_j$&$b_j$&$a_j$&$b_j$&$a_j$&$b_j$\\\midrule
0 & 1 & 1 & 1 & 1 & 1 & 1 \\
1 & 0 & 0 & 0 & 0 & 0 & 0 \\
2 & 0 & 0 & 0 & 0 & 0 & 0 \\
3 & 0 & 0 & 0 & 0 & 0 & 0 \\
4 & 0 & 0 & 0 & 0 & 0 & 0 \\
5 & 0 & 0 & 0 & 0 & 0 & 0 \\
6 & 4 & 588 & 0 & 564 & 0 & 564 \\
7 & 0 & 2144 & 16 & 2240 & 12 & 2208 \\
8 & 130 & 8538 & 102 & 8298 & 114 & 8394 \\
9 & 480 & 28960 & 504 & 29312 & 508 & 29344 \\
10 & 1436 & 79436 & 1436 & 79396 & 1392 & 79044 \\
11 & 2688 & 171968 & 2688 & 171648 & 2712 & 171840 \\
12 & 4312 & 300808 & 4288 & 300776 & 4344 & 301224 \\
13 & 6464 & 417344 & 6416 & 417280 & 6360 & 416832 \\
14 & 7260 & 445700 & 7416 & 446748 & 7392 & 446556 \\
15 & 5760 & 357856 & 5680 & 357056 & 5724 & 357408 \\
16 & 2981 & 201501 & 2905 & 201005 & 2901 & 200973 \\
17 & 992 & 70304 & 1080 & 71040 & 1068 & 70944 \\
18 & 260 & 12004 & 236 & 11788 & 240 & 11820 \\
\bottomrule
\end{tabular}
\caption{Complete stabilizer and normalizer weight distributions. Logical-label coefficients are $\ell_j=b_j-a_j$.}\label{tab:enumerators}
\end{table}

The coefficient sums are
\[
 \sum_{j=0}^{18}a_j=32768=2^{15},\qquad
 \sum_{j=0}^{18}b_j=2097152=2^{21}
\]
for all three codes. For $\CA$, the $588$ weight-six normalizer labels consist of $4$ stabilizers and $584$ labels in $N_A\setminus S_A$. For each derivative, all $564$ weight-six normalizer labels are outside its stabilizer.

As an independent enumerator consistency test, the quaternary MacWilliams identity is
\begin{equation}\label{eq:macwilliams}
 B(u,v)=\frac1{|S|}A(u+3v,u-v).
\end{equation}
In coefficient form, for every $0\le j\le18$,
\begin{equation}\label{eq:krawtchouk}
 2^{15}b_j=\sum_{i=0}^{18}a_iK_j(i),\qquad
 K_j(i)=\sum_{h=0}^j(-1)^h3^{j-h}\binom{i}{h}\binom{18-i}{j-h},
\end{equation}
where out-of-range binomial coefficients are zero. All $19$ identities are verified with exact integers for each code. This check corroborates the enumeration; it is not substituted for a computation of the physical normalizer or for the witness outside the stabilizer.

\section{Pairwise local-unitary inequivalence}\label{sec:inequivalence}
\begin{theorem}\label{thm:inequivalence}
The codes $\CA,\CB,\CC$ are pairwise inequivalent under products of arbitrary single-qubit unitaries and permutations of the $18$ qubits.
\end{theorem}
\begin{proof}
For an orthogonal code projector $\Pi$ of rank $K$, define
\begin{equation}\label{eq:projectorinvariant}
 \mathsf A_j(\Pi)=\frac1{K^2}\sum_{\substack{P\in\{I,X,Y,Z\}^{\otimes18}\\\wt(P)=j}}
 |\Tr(P\Pi)|^2.
\end{equation}
For a stabilizer code, $\Pi=|\mathcal S|^{-1}\sum_{g\in\mathcal S}g$. Pauli trace orthogonality implies that $|\Tr(P\Pi)|=K$ when the binary label of $P$ belongs to $S$, and zero otherwise. Signs of the actual stabilizer elements do not change the squared traces. Consequently $\mathsf A_j(\Pi)=a_j$.

Conjugation by a single-qubit unitary fixes $I$ and acts as a real orthogonal transformation on the real span of $X,Y,Z$. For any fixed support $T\subseteq\{1,\ldots,18\}$, a product unitary therefore acts orthogonally on the vector of Pauli trace coefficients having support exactly $T$. Their squared norm is unchanged. Summing this norm over supports of size $j$ proves invariance of~\eqref{eq:projectorinvariant} under arbitrary product unitaries. A qubit permutation only permutes the supports and also preserves it.

The coefficient $a_6$ equals $4$ for $\CA$ and $0$ for both derivatives, distinguishing $\CA$ from each of them. The coefficient $a_7$ equals $16$ for $\CB$ and $12$ for $\CC$, distinguishing those two. Thus no product unitary followed or preceded by a qubit permutation maps any of the three code spaces onto another.
\end{proof}

This is stronger than inequivalence under local Clifford operations and permutations. It does not assert inequivalence under general entangling Clifford operations: equation~\eqref{eq:exactclifford} explicitly supplies such equivalences. Nor does the argument classify all pure $[[18,3,6]]$ codes or prove uniqueness of any one enumerator.

\section{Reproducible finite certificates}\label{sec:reproducibility}
\subsection{Primary portable verification}
The accompanying directory \texttt{verification/} contains a standard-library-only Python certificate. From the manuscript directory, run
\begin{center}\ttfamily python3 verification/verify\_all.py\end{center}
with Python $3.10$ or newer and without optimization flags. Assertions are part of the certificate; the script explicitly refuses to run with \texttt{-O}. No internet connection, numerical solver, search database, compiler, third-party Python package, or precomputed normalizer basis is required.

For each code, the individual self-contained verifier embeds the $15$ physical generator labels. It performs the following exact steps.
\begin{enumerate}
\item Check that the labels lie in $V_{18}$, have binary rank $15$, and are pairwise symplectically orthogonal.
\item Convert each normalizer constraint $[s,r]=0$ into an ordinary binary dot-product equation by swapping the $x$ and $z$ halves of $s$. Binary row reduction gives a fresh nullspace basis of rank $21$.
\item Enumerate the $2^{15}$ stabilizer labels and store them in a set. Enumerate the $2^{21}$ normalizer labels, tally the physical support weight of each, and test set membership to count labels outside the stabilizer.
\item Check the exact minimum nonzero normalizer weight, the logical minimum, all $19$ MacWilliams coefficients, the coefficient sums, and the stated weight-six witness by direct weight, commutation, and nonmembership tests.
\end{enumerate}
Enumeration uses a binary-reflected Gray order: from label number $t-1$ to $t$, toggle basis vector $1+\nu_2(t)$. The Gray sequence visits every coefficient vector once; independence of the basis ensures that every physical label is visited exactly once. Therefore the finite exhaustive calculation proves absence of low-weight labels, rather than sampling their absence.

The combined verifier also recomputes all $1024$ local coset minima on each side, checks every nonzero quotient inequality, verifies the quotient ranks and dual forms, checks~\eqref{eq:poisson}, and reconstructs the exact generator list of $\CA$. It then verifies both transvections on all $36$ standard basis vectors, all pairings of those vectors, the displayed Pauli strings, and every signed generator conjugation. Expected output is preserved in \texttt{verification/EXPECTED\_VERIFICATION.txt}; it is a comparison aid, not an input to the calculation of the code parameters.

\subsection{Independent low-weight physical error checks}
Each code directory includes a short C++17 verifier with its physical generators embedded. It does not use the normalizer basis or the local gluing construction. For each support of size $w=1,\ldots,5$, it visits every one of the $3^w$ nonidentity Pauli assignments and computes its $15$-bit syndrome. The number tested per code is exactly
\begin{equation}
 \sum_{w=1}^{5}\binom{18}{w}3^w
 =54+1377+22032+247860+2082024
 =2353347.
\end{equation}
All these syndromes are nonzero. Hence no nonidentity Pauli of weight at most five belongs to the normalizer, independently verifying the lower distance bound. The full normalizer enumeration and the explicit logical witness supply the upper bound.

For example, compile and run the check for $\CA$ as follows, replacing \texttt{A} by \texttt{B} or \texttt{C} for the other codes:
\begin{center}\small\ttfamily
c++ -O3 -std=c++17 verification/frozen\_A/check\_errors.cpp -o check\_A\\
./check\_A
\end{center}
The integer arithmetic in both implementations is exact. The checks deliberately use two complementary traversals: basis combinations for the full normalizer and support-by-support physical errors for the low-weight exclusion. They are separately implemented checks.

\subsection{Construction versus search}
The finite data in Appendix~\ref{app:local}, the columns of $F$, and the two transvection labels completely specify the constructions. Verification does not require reproducing a random search or trusting its stopping condition. The gluing proof explains why the base normalizer has minimum weight six; the derivative certificates check the property again because a nonlocal Clifford can alter weight. Search histories and selected enumerator fingerprints are not used as substitutes for the generators or the exact certificate.

\section{Conclusion and scope}
The three concrete codes establish a small family of pure $[[18,3,6]]$ stabilizer spaces with distinct product-local-unitary equivalence types. The first combines two local coset metrics through an isometry of degenerate dual alternating forms. Two specified four-site transvections change its low-weight stabilizer spectrum while preserving distance six, as verified independently in physical coordinates. The complete generator, phase, enumerator, and witness data make the construction usable without access to a discovery pipeline.

The distinction between global and local equivalence is essential: a simple global Clifford relationship coexists with a rigorous obstruction to every product-unitary and permutation equivalence. The next questions include whether additional transvections give further local equivalence classes, how systematically the quotient method can generate short codes, and whether stronger length-$18$ parameters can be achieved. No completeness claim about the transvection orbit, optimal encoded dimension, or the $[[18,4,6]]$ parameter set is made here.

\appendix
\section{All physical stabilizer generators}\label{app:generators}
Every row in Table~\ref{tab:generators} is assigned positive Hermitian phase. The row order is used in the signed transformations in Section~\ref{sec:derivatives}. Pauli words run from qubit $1$ on the left to qubit $18$ on the right, consistently with~\eqref{eq:int}. The binary span of the labels is $S_D$ for code $\mathcal Q_D$; a product of displayed positive generators can itself have a negative Hermitian phase, which does not alter its binary label or weight.

\begin{table}[htbp]
\centering\small
\setlength{\tabcolsep}{6pt}
\begin{tabular}{@{}rlll@{}}\toprule
$j$&$\CA$&$\CB$&$\CC$\\\midrule
1 & \texttt{XIIIZIIYZZZIIIIYZI} & \texttt{ZIIIZIIIZZZIIXYYZI} & \texttt{XIIIZIIYZZZIIIIYZI} \\
2 & \texttt{IYIZIIZYZIZIZZZIZY} & \texttt{IYIZIIZYZIZIZZZIZY} & \texttt{IYIZIIZYZIZIZZZIZY} \\
3 & \texttt{IIXIZIZXIIZIYIIXZY} & \texttt{YIXIZIZZIIZIYXYXZY} & \texttt{IYXIZIXXIIZIZIIZZY} \\
4 & \texttt{IZZYZZIYIZZZIZIIII} & \texttt{YZZYZZIIIZZZIYYIII} & \texttt{IXZYZZYYIZZZXZIYII} \\
5 & \texttt{IIIIYZIXZIZZYIIYII} & \texttt{YIIIYZIZZIZZYXYYII} & \texttt{IYIIYZYXZIZZZIIIII} \\
6 & \texttt{IZZIZXZXIIIIXZIZZX} & \texttt{IZZIZXZXIIIIXZIZZX} & \texttt{IXZIZXXXIIIIIZIXZX} \\
7 & \texttt{IZIIZIYYZIZIYZZYIX} & \texttt{IZIIZIYYZIZIYZZYIX} & \texttt{IZIIZIYYZIZIYZZYIX} \\
8 & \texttt{IIZZIIZZYZZIYZIYZZ} & \texttt{IIZZIIZZYZZIYZIYZZ} & \texttt{IIZZIIZZYZZIYZIYZZ} \\
9 & \texttt{IIIZZZZIIXIIIIZXZI} & \texttt{YIIZZZZYIXIIIXXXZI} & \texttt{IIIZZZZIIXIIIIZXZI} \\
10 & \texttt{IZZIZIIZIZYZZIZYZZ} & \texttt{IZZIZIIZIZYZZIZYZZ} & \texttt{IZZIZIIZIZYZZIZYZZ} \\
11 & \texttt{IIZZZZIIIZZYIZZYII} & \texttt{IIZZZZIIIZZYIZZYII} & \texttt{IIZZZZIIIZZYIZZYII} \\
12 & \texttt{IZZZZIIIIIIZZYIIZI} & \texttt{YZZZZIIYIIIZZZYIZI} & \texttt{IZZZZIIIIIIZZYIIZI} \\
13 & \texttt{IZIZIZIZZZIIIIXXZZ} & \texttt{IZIZIZIZZZIIIIXXZZ} & \texttt{IZIZIZIZZZIIIIXXZZ} \\
14 & \texttt{IIIIIIIIIIZZZZZZXX} & \texttt{IIIIIIIIIIZZZZZZXX} & \texttt{IIIIIIIIIIZZZZZZXX} \\
15 & \texttt{ZZZZZZZZIIIIIIIIII} & \texttt{ZZZZZZZZIIIIIIIIII} & \texttt{ZZZZZZZZIIIIIIIIII} \\
\bottomrule
\end{tabular}
\caption{The complete $15$-row generator bases. All words have length $18$.}\label{tab:generators}
\end{table}

\section{Complete local gluing data}\label{app:local}
All entries in this appendix are nonnegative integers using~\eqref{eq:int} with the \emph{local} length, $n=8$ on the left and $n=10$ on the right. The local blocks are embedded into global qubits $1,\ldots,8$ and $9,\ldots,18$, respectively. Explicitly, for a local integer label $r=x+2^n z$ and a zero-based block offset $o\in\{0,8\}$,
\begin{equation}
 \iota_{n,o}(r)=2^o x+2^{18+o}z.
\end{equation}
Thus the $21$ spanning labels for $N_A$ are the four embedded $C$ rows, the seven embedded $E$ rows, and
\[
 \iota_{8,0}(q_L(e_i))+\iota_{10,8}(q_R(Fe_i)),\qquad i=1,\ldots,10,
\]
where addition means XOR.

\begin{center}
\begin{tabular}{@{}crr@{}}\toprule
Row&$C$ on eight qubits&$E$ on ten qubits\\\midrule
1&65280&209139\\
2&255&52428\\
3&3891&258816\\
4&13141&4557\\
5& &25930\\
6& &74594\\
7& &786483\\\bottomrule
\end{tabular}
\end{center}
The marked subspace $A$ is spanned by rows $1,2$ of $C$; the marked subspace $B$ is spanned by rows $1,2,3$ of $E$.

\begin{center}
\begin{tabular}{@{}crrrr@{}}\toprule
$i$&$q_L(e_i)$&$q_R(e_i)$&$d_{L,i}$&$d_{R,i}$\\\midrule
1&41092&663560&24702&812934\\
2&24705&1040385&20552&594171\\
3&40993&188441&24596&91366\\
4&24596&245763&20498&70425\\
5&33153&966736&3&619650\\
6&3&589829&16920&865104\\
7&33952&499978&12&88022\\
8&12&51&18442&283146\\
9&129&24600&13141&88037\\
10&132&115024&15462&29831\\\bottomrule
\end{tabular}
\end{center}
The first two columns specify the quotient lifts. The last two are convenient dual representatives satisfying
\[
 [d_{L,i},C]=0,\quad[d_{L,i},q_L(e_j)]=\delta_{ij},\qquad
 [d_{R,i},E]=0,\quad[d_{R,i},q_R(e_j)]=\delta_{ij}.
\]
Their pairwise Gram matrices are precisely~\eqref{eq:J}. They need not lie in the spans of the corresponding primal lifts, and the Gram matrix of the primal lift vectors is not the form that $F$ must preserve.

For a ten-bit label $t=\sum_i t_i e_i$, compute $q_L(t)=\sum_i t_iq_L(e_i)$ and $Ft=\sum_i t_iFe_i$ by XOR. The local cost tables are thereby unambiguously determined by the displayed finite data; no extra search-specific convention is needed. The certificate recomputes both tables from these definitions and compares them with the supplied JSON copies.

\section{Certificate file map}\label{app:files}
The manuscript bundle contains:
\begin{itemize}
\item \texttt{manuscript.tex}, \texttt{references.bib}, and the compilation instructions: the paper source and references. A compiled PDF is a subsequent typesetting deliverable.
\item \texttt{verification/verify\_all.py}: the combined exact verifier for the constructions, phases, and enumerations.
\item \texttt{verification/frozen\_A/}, \texttt{frozen\_B/}, and \texttt{frozen\_C/}: standalone parameter verifiers, explicit rows and Pauli words, expected enumerators, and direct low-weight C++ checks. These are relative subdirectories of \texttt{verification/}.
\item \texttt{verification/construction/}: the local quotient data and independent gluing reconstruction, together with \texttt{verification/RAW15.rows} as its expected output basis.
\item \texttt{verification/enumerators.csv}: all $a_j,b_j,\ell_j$ for all three codes.
\item \texttt{verification/audit\_published\_example.py}: the standalone check of the earlier final published matrix in Appendix~\ref{app:prior}.
\item \texttt{README.md}: portable verification and build instructions; \texttt{SHA256SUMS}: final integrity hashes.
\end{itemize}
The expected-result JSON files and logs are not sufficient evidence alone: the scripts recompute the relevant quantities from physical generator data. Hashes detect accidental changes, while the mathematical checks establish the claimed properties of those data.


\section{Check of the earlier published example}\label{app:prior}
The final article of Cohen, Encheva, and Litsyn \cite{CohenEnchevaLitsyn1999}, page 2496, gives the following twelve rows for the binary extension code $C'$ in Example 1. They agree with the rows in the inspected 1998 preprint. This appendix audits those displayed data, not every possible construction with the stated parameters.
\begin{center}\small\ttfamily
\begin{tabular}{@{}ll@{}}
111111000000000000 & 001000001100101010\\
001011100000101000 & 001000001100010101\\
001001011000110000 & 011000000111000100\\
011100000100110000 & 110000101000000000\\
011010000000011100 & 000000001000110001\\
011001001010001000 & 000000010100001001
\end{tabular}
\end{center}
The first six rows are in the left column and the remaining six in the right column. These are classical binary rows, with the ordinary Euclidean pairing, not symplectic Pauli labels.

Let $M$ be this $12\times18$ matrix. Binary elimination gives
\[
 \rank M=12,\qquad \rank(MM^{\mathsf T})=10.
\]
Every row is even, so the Euclidean pairing restricted to $C'$ is alternating. Its radical therefore has dimension $12-10=2$, and its nondegenerate quotient is a $10$-dimensional symplectic space. Any totally isotropic subspace of $C'$ has dimension at most
\[
 2+\frac{10}{2}=7.
\]
The cited theorem requires a length-$18$, dimension-$9$ code $C$ with $C^\perp\subseteq C\subset C'$. Since $\dim C^\perp=9$, this would force $C=C^\perp$, making $C$ a $9$-dimensional totally isotropic subspace of $C'$. The bound above rules this out for the displayed matrix, independently of how one might choose nine rows or a different basis for $C$.

The portable script \texttt{verification/audit\_published\_example.py} embeds the twelve displayed rows and checks the ranks and dimension bound exactly. A final-page inspection verifies the correspondence with the published data. The conclusion is restricted to the displayed construction's failure to meet the theorem hypothesis. It neither disproves $[[18,3,6]]$ existence nor establishes the absence of later corrections or other valid constructions.

\section*{Acknowledgments}
Claude helped with the search, the proofs, the verification code, and the
manuscript.

\bibliographystyle{plain}
\bibliography{references}
\end{document}
